\section{Methodology}
\label{sec:method}

\subsection{Task Decomposition and Planning}
\label{sec:planning}

The \texttt{parse\_intent} node asks a lightweight model for a structured intent: an action verb, a target site, optional target and content strings, a risk level and a confidence. If the request contains a connective (``and'', ``then'', ``after'', a comma or semicolon) or more than eight words, a second call asks a planning model for an ordered plan $P=(g_1,\dots,g_n)$ in which every step carries a description, an action type and an expected outcome. When the request asks for literal text to be entered, the plan is replaced by a single text-entry step. Otherwise a one-step plan is created from the intent.

Planning then proceeds one action at a time. For the current step, \texttt{plan\_action} builds a prompt from the goal, the plan with the current step marked, the page title, URL and state, and the list of interactive elements, and asks for one action from a fixed vocabulary (\texttt{navigate}, \texttt{click}, \texttt{type\_text}, \texttt{press\_key}, \texttt{scroll}, \texttt{extract}, \texttt{complete}, \texttt{need\_help}). The model's choice is then passed through deterministic overrides: exact-text requests are grounded to the highest-scoring visible text input; repeated navigation to the current site is suppressed; and plan steps whose descriptions name clicking a result, extracting, or reporting are mapped to the corresponding action type. Some overrides are site-specific. On a Google page with a search goal the agent types the query and presses Enter, and a planned navigation straight to a Google results URL is rewritten into typing into the search box. We discuss the ethics of this choice in Section~\ref{sec:security}.

\subsection{Grounding: DOM First, Vision on Demand}
\label{sec:grounding}

Each observation starts with a DOM pass. The extractor queries buttons, inputs, text areas, selects, links, elements with button or link roles, and editable regions. For each it records tag, role, accessible label, the first 120 characters of text or value, placeholder, input type, visibility, whether it is enabled, a bounding box, and three selectors: an attribute selector on an identifier that the extractor writes into the live DOM, a CSS selector, and an XPath. Only interactable elements are kept. They are ranked with submit-like controls first, then text fields, selects and links; duplicates are removed, links are capped at five, cookie and advertising labels are dropped, and the manifest is truncated to 50 elements. The executor resolves an element by trying the three selectors in order.

The vision model is invoked only when the planner returns \texttt{need\_help}, either by choice or because its call failed. The agent then takes a screenshot and asks the vision model for an action and a target point expressed as fractions of the viewport. The point is scaled to the 1280\,$\times$\,800 viewport and clicked, and any text is typed at that position. A vision action is recorded as dispatched without a follow-up check; the next pass of the completion predicate is what catches a wrong click. A cache keyed on a hash of the screenshot, goal and target avoids repeated inference on identical inputs.

% Formal model of evidence-gated completion (Sec. IV-C), as implemented after the fixes (commit eb3803e).
\subsection{Evidence-Gated Verification Model}
\label{sec:formal}

We model an episode as a sequence of browser states $s_t$ with observations $o_t=\Omega(s_t)$ (URL, title, element manifest, screenshot). The planner maps the goal $g$, $o_t$ and retrieved context to an action $a_t$; the executor applies it and returns a dispatch flag $\alpha_t\in\{0,1\}$, which is $1$ when the browser call raised no exception and, for text entry, the value read back matched. Some fallback paths (a mouse-and-keyboard retry of text entry, vision-coordinate clicks) return $\alpha_t=1$ without such a check.

\paragraph{Blocking predicate} $\beta(s)=1$ when the page is a bot check: the URL contains a challenge path, the title contains challenge wording (\emph{captcha}, \emph{not a robot}, \emph{sorry\ldots}), the body contains one of 12 challenge phrases, or a visible challenge frame other than an invisible reCAPTCHA is present. This check runs before $\Phi$ and ends the episode as \textsc{Blocked}.

\paragraph{Step and task predicates} The decomposer produces a plan $P=(g_1,\dots,g_n)$. The per-step predicate is still weak, $\phi_i := \alpha_t$: a step is marked complete when an action other than \texttt{complete} dispatches during it. Completion is decided by $\Phi(g,s_{t+1},\hat{y},P)$, evaluated on the live page after the action:
\begin{equation}
\Phi = \neg E \wedge \neg H \wedge Q \wedge T \wedge X \wedge R,
\label{eq:phi}
\end{equation}
where $E$ holds for a browser-error URL, an error or ``not found'' title, or a redirect to a sign-in page that the request did not name; $H$ holds when a search or extraction request, or a multi-step plan, is still on the bare home page of one of four search engines; $Q$ requires a results page to carry the query quoted in $g$; $T$ requires the literal text that $g$ asks to be entered to appear in an input, text area or editable element, exactly if $g$ says ``exactly''; $X$, for extraction requests, requires an answer or structured extraction that mentions at least one content word of $g$ and does not merely restate a navigation step; and $R$, for multi-step plans, requires that no step is pending unless a structured extraction exists. $H$, $Q$, $T$, $X$ and $R$ are instantiated only when their triggers occur; when none does, $\Phi$ reduces to $\neg E$.

\paragraph{When $\Phi$ is evaluated} $\Phi$ is evaluated when the planner emits \texttt{complete}, when the plan index has reached $n$, or when a requested text entry has dispatched. For a one-step plan, the default for short requests, this is after every successful action, so $\Phi$ also acts as the termination rule. The answer $\hat{y}$ is the text an extraction grounded in page lines, or else the planner's summary on \texttt{complete}; the reasoning attached to ordinary actions is not an answer.

\paragraph{Completion rule} The agent reports success only if $\Phi=1$; each evaluation is written to the evidence store with its expected and observed values. A passing evaluation can still yield \textsc{Waiting approval}, and the terminal node downgrades a success to \textsc{Failed} if the intent named a site but no navigation was verified. A rejected evaluation does not end the episode: control returns to observation and planning until a counter of intent, decomposition and planning calls reaches $K=15$.

\paragraph{Consequence} Dispatch success does not imply reported success, $\alpha_t = 1 \not\Rightarrow \Phi=1$; a successful \texttt{navigate} to a search engine's home page for a search request is an example ($H=1$). The converse, that $\Phi=1$ implies the goal is met, does not hold, and Section~\ref{sec:results} measures how often it fails.

\paragraph{False completion rate} Let $\mathcal{T}_c$ be the episodes reported complete and $\Phi^{\star}$ an independent, task-specific oracle written by the experimenter. Then
\begin{equation}
\mathrm{FCR} = \frac{\lvert\{\tau\in\mathcal{T}_c : \Phi^{\star}(\tau)=0\}\rvert}{\lvert\mathcal{T}_c\rvert}.
\label{eq:fcr}
\end{equation}
FCR is a precision-type measure; an agent that never completes has no false completions. It must therefore be reported together with task success and the false-rejection rate, the fraction of goal-satisfying episodes not reported complete.


\subsection{Failure Classification and Recovery}
\label{sec:recovery}

When an action fails, the \texttt{verify} node first retries it through a fresh observation, up to three times. Further failures reach the recovery node, which classifies the error message by pattern into one of seven types (\texttt{captcha}, \texttt{transient}, \texttt{element\_not\_found}, \texttt{verification\_failed}, \texttt{navigation\_failed}, \texttt{vision\_needed}, \texttt{llm\_failure}) or \texttt{unknown}, and records a diagnosis. A CAPTCHA ends the task as \textsc{Blocked}. Other failures are assigned a strategy from the ladder \texttt{retry} $\rightarrow$ \texttt{alternative\_selector} $\rightarrow$ \texttt{vision\_fallback} $\rightarrow$ \texttt{replan}, advancing one level every two attempts, with type-specific shortcuts (for example, \texttt{vision\_needed} maps directly to \texttt{vision\_fallback}); after six attempts the task fails. Every decision is stored as a recovery record in the evidence directory. In the current implementation the selected strategy is recorded but not acted on: the state schema does not declare the strategy field, so it is dropped, and every non-terminal recovery returns to observation and replanning. Differentiated strategy execution is therefore part of the design but not of the running system (Table~\ref{tab:status}).

\subsection{Role-Based Local Model Routing}
\label{sec:routing}

Small local models differ in what they do well, and loading a model has a cost. The router therefore assigns each call a role and picks a model for that role. A deterministic analyzer maps the call's type, text and flags to required capabilities and a role: recovery calls go to a recovery model; calls with an image, or with screen-related keywords, go to a vision model; intent parsing and simple navigation go to a lightweight model; planning, multi-clause or long requests go to a planner; everything else goes to a general model. Candidates for each role come from a configured list (Table~\ref{tab:routing}) filtered by what the local Ollama server reports as installed, with the list refreshed every 300\,s. To avoid reloading models, the router keeps the currently active model whenever it already has the required capabilities, except on recovery calls. If no candidate is available the router falls back to a configured fallback model and then to the base model. A single-model mode, used for ablation, sends all text calls to one model. A dedicated cool-down parameter exists in the configuration but is not read by the router.

\begin{table}[t]
\centering
\caption{Default role configuration of the model router (Ollama tags, in preference order). The router uses the first candidate that is installed. The launcher installs only \texttt{qwen2.5:1.5b} and one vision model by default, in which case every text role resolves to \texttt{qwen2.5:1.5b}.}
\label{tab:routing}
\footnotesize
\setlength{\tabcolsep}{3pt}
\begin{tabular}{@{}lp{3.05cm}p{2.85cm}@{}}
\toprule
Role & Selected when & Candidates \\
\midrule
Lightweight & intent parsing; simple navigation & \texttt{deepseek-r1:1.5b}, \texttt{qwen3.5:2b}, \texttt{qwen2.5:1.5b} \\
Planner & decomposition; multi-clause or long requests & \texttt{qwen3.5:2b}, \texttt{qwen2.5:7b}, \texttt{deepseek-r1:1.5b} \\
Vision & image input or screen keywords & \texttt{qwen3-vl:2b}, \texttt{moondream} \\
Coder & coding keywords & \texttt{qwen2.5-coder:3b}, \texttt{qwen2.5-coder:1.5b} \\
Recovery & error recovery & \texttt{qwen2.5:7b}, \texttt{qwen3.5:2b} \\
General & otherwise & \texttt{qwen2.5:1.5b}, \texttt{qwen2.5:7b} \\
Fallback & no candidate installed & \texttt{qwen2.5:1.5b} \\
\bottomrule
\end{tabular}
\end{table}


\subsection{Episodic Memory}
\label{sec:memory}

At the end of each task, \texttt{complete} writes a record of the outcome to a SQLite table and to a Chroma collection: on success, a strategy note with the goal, the final URL and the number of model calls; on failure, the error and URL. The runner also stores a one-line episode summary. The retrieval side is built: a context router decides per task whether to query a knowledge collection and always queries memory, a context builder deduplicates the results and fits them into a 3{,}500-token budget split 65/35 between knowledge and memory, and the result is placed in the planning prompt under an explicit untrusted-content banner. On the live execution path, however, the runner initializes both retrieved lists as empty rather than unset, and the router interprets this as ``already retrieved''. Memory is therefore written but not read during normal runs. Because the stored strategy holds no selectors or action sequences, fixing the read path alone would give the planner little to reuse; we treat evidence-linked procedural memory as future work.

\subsection{Extension to Local System Access}
\label{sec:desktop}

The browser loop generalizes to the desktop if each web-specific part has a local counterpart: the accessibility tree plays the role of the DOM, a window screenshot the role of the page screenshot, and typed operating-system predicates the role of URL and DOM checks. We first state what exists. The repository contains a desktop executor that wraps PyAutoGUI for clicking, typing, hotkeys and screenshots and psutil for listing processes. No graph node, action type or API route calls it, and its only test checks that its methods exist; without a display it degrades silently (a fixed screen size and empty screenshots). A file helper can read and write text files but is not called by the agent and does not confine paths to its workspace. A file-existence predicate exists in the verification manager and is unit-tested but unused. Everything below is a design (\status{Proposed}).

\paragraph{Action space and grounding} We define file actions (list, stat, hash, read, create, overwrite with backup, copy, move, move to trash, permanent delete), process actions (list, launch an allow-listed executable, terminate), application actions (open, focus, close) and UI actions (invoke, set value, type, hotkey, clipboard). Grounding reads the accessibility tree of the foreground window through Windows UI Automation, the macOS Accessibility API or AT-SPI on Linux, and produces a manifest with the same shape and element budget as the DOM manifest. As on the web, element names pass through the sanitizer. Vision grounding on a screenshot of the target window is used when the tree exposes too few actionable elements (canvas or GPU-rendered applications, some Electron applications) or when the planner asks for help; unlike the current web path, a coordinate action counts as successful only after its postcondition holds.

\paragraph{Desktop verification predicates} Each mutating action carries typed, exact postconditions that are evaluated after execution and stored with before/after values (Table~\ref{tab:desktop}). File predicates compare SHA-256 digests, sizes and modification times; process predicates identify a process by PID and creation time, which guards against PID reuse; window and accessibility predicates compare titles, focus, element values and states. Read-only actions are checked for absence of side effects by comparing a manifest of the workspace before and after. The task-level rule is unchanged: completion requires the predicates of every step and the goal.

\paragraph{Permission tiers} The tier of an action is computed deterministically from the concrete action and its arguments, never from the model's risk label. \emph{T0} (observe) actions are read-only and run without approval. \emph{T1} (reversible) actions have a recorded inverse, such as a new file, an overwrite whose pre-image is backed up, a move, or a move to trash; they run automatically only if the user has enabled automatic approval of low-risk actions. \emph{T2} (destructive) actions have no reliable inverse or affect state the agent did not create, such as permanent deletion, terminating a process the agent did not launch, or invoking a control whose label matches a destructive lexicon (Delete, Discard, Send, Pay); they always need approval. \emph{T3} actions can move data off the machine or touch secrets, such as launching network-capable programs or typing into password fields; they always need approval, and secrets are injected by the executor without passing through the model. Paths outside allow-listed roots, system and credential directories, shell interpreters with inline code, and privilege elevation are denied outright.

\paragraph{Enforcement and audit} A new \texttt{action\_gate} node between planning and execution computes the tier, applies the path jail (resolving symbolic links before the containment test), produces a dry-run preview, and either allows the action, pauses for approval, or rejects it. An approval is bound to a hash of the exact action so that the model cannot change arguments after the user approves. Processes are launched without a shell, with a scrubbed environment and resource limits. Each local action is appended to the audit log with digests of its arguments and of the before/after state, never with file contents or typed text, and each log line carries the hash of the previous line so that tampering is detectable. The same gate can serve web actions such as form submission, which today are covered only by the task-level approval.

\begin{table}[t]
\centering
\caption{Proposed desktop postconditions and default permission tiers (design, not implemented). $h_0$, $t_0$: values captured before the action; $(\mathit{pid}, t_c)$: process ID and creation time.}
\label{tab:desktop}
\footnotesize
\setlength{\tabcolsep}{2.5pt}
\begin{tabular}{@{}p{1.75cm}p{4.35cm}c@{}}
\toprule
Action & Postcondition checked after execution & Tier \\
\midrule
file read, list, hash & workspace manifest unchanged & T0 \\
create file & file exists $\wedge$ $\mathrm{sha256}=\mathrm{sha256}(\text{content})$ & T1 \\
overwrite & backup of pre-image stored $\wedge$ hash equals new content & T1 \\
move & source absent $\wedge$ $\mathrm{sha256}(\text{dst}) = h_0(\text{src})$ & T1 \\
move to trash & path absent $\wedge$ trash entry with that origin & T1 \\
permanent delete & path absent & T2 \\
launch app & $\mathit{pid}$ alive with creation time $t_c$ $\wedge$ window title matches & T1/T3 \\
terminate & $(\mathit{pid}, t_c)$ gone within timeout & T1/T2 \\
focus window & foreground title matches pattern & T0 \\
set value & accessibility value of element equals input & T1/T3 \\
toggle control & element state contains \textit{checked} & T1/T2 \\
\bottomrule
\end{tabular}
\end{table}

